\section{How We Tested}
\label{sec:setup}

\subsection{Measuring and Comparing Accuracy}
\label{sec:stats}

In every time step, an exact search over all stored passages finds each test query's true \RNumPlainTopK{} best matches and when each was added.
\emph{Recent recall} is the share of a query's true best matches from the current and previous time steps that the system returns.
Only queries with at least one recent match count toward it: about \RNumPlainQueriesWithRecentPctSteady{} of the mixed-topic queries and \RNumPlainQueriesWithRecentPctNewTopic{} of the new-topic queries.
\emph{Overall recall} counts all \RNumPlainTopK{} true best matches, old or new.
We give recall as a percentage and call a difference of one percentage point one \emph{point}.

We run each main setting five times, with seeds numbered \RNumSeedMin{} to \RNumSeedMax{}.
A seed is a number that fixes a run's random choices, the way a recorded shuffle fixes the order of a deck of cards.
It sets the arrival order, the clusters, the queries asked, and the shards' search graphs.
The passages are the same in every run, and all methods see the same queries, so we compare methods query by query.

To compare two methods, we average the difference in recall between them over all queries and seeds.
Other seeds or queries would give a somewhat different average, so we ask how far it could move by chance.
Our main test is a \RNumCIPct{} range found by resampling.
We redraw the five seeds with repeats allowed, like drawing names from a hat and putting each one back.
Within each drawn seed, we redraw its queries the same way.
We recompute the average difference for each of \RNumPlainBootReps{} such redraws, and the middle \RNumCIPct{} of these averages is the range.
If the whole range lies on one side of zero, chance alone would rarely produce the difference, and we call it \emph{clear}.
Otherwise we say there is no clear difference.

Our simpler test uses only the five seed averages of the difference.
Their \emph{spread} is their standard deviation, roughly the typical distance of one seed's average from the mean of all five.
This textbook range is centered on that mean, and its half-width is about \RNumPlainTCritFive{} times the spread divided by the square root of five.
The large factor allows for having only five numbers.
The two tests can disagree because resampling also counts how single queries vary, but with only five seeds it understates how much seeds differ.
So when seeds differ a lot, as in several main eight-shard tests in Table~\ref{tab:main}, our range was \RNumBootNarrowerMin{} to \RNumBootNarrowerMax{} narrower.
When the five seed averages agree closely, the simpler range shrinks and ours can be wider: \RNumBootWiderMThreeRNMinusQLNewTopic{} wider for the probe against most-searched (Section~\ref{sec:res-which}).

A clear lead is \emph{fragile} if it fails any one of four further checks.
Its range must not nearly touch zero, and the simpler test must exclude zero.
No seed may point the other way, and a most-changed lead must survive paying for its score messages (Section~\ref{sec:follow}).

Our ranges treat each asked query as a separate, unrelated sample, although the same \RNumPlainTestQueries{} queries repeat about \RNumQueryRepeatsQTwoHundred{} times per run.
They also rest on the one sample of passages and queries we drew from the data set.
Nor are they widened for our many comparisons, which raises the chance that some clear result is a fluke.
Against other options we use the \emph{best refresh rule}: whichever of the five rules has the highest recent recall in that setting.
Picking it from the same runs we compare favors refresh a little.
This makes the alternatives' tests harder to pass under starting-center placement, and refresh's tests slightly easier under eight-shard placement.

\subsection{Our Plan, Written Before the Main Analysis}

Writing tests down before seeing the results keeps researchers from reporting only the tests that happen to look good.
Before any run, we stored a plan with first settings, message sizes, and three hypotheses, or predictions.
One was that the probe would beat the other rules by at least \RNumPlainGateThreshold{} points under starting-center placement.
A small trial run, or pilot, with one seed (\RNumPilotSeed{}) then showed that refresh gained little there, and all three were replaced.

Before the main runs, a code review and a first written change to the plan, called an addendum, revised the design.
Together they added the new-topic stream, one-shard placement, the three query rates, and an ID number in each pushed passage.
They also replaced a rule that slowly raised waiting shards' scores with a waiting limit set by the budget.
The first addendum set five of the tests below.
Two of them restate pilot results and reuse the pilot's seed, so Section~\ref{sec:res-near} also checks them on the other four seeds.
Testing our analysis program on a partial set of \RNumPartialMOneTestedRuns{} starting-center runs printed one result, which changed no test or rule.
A second addendum came when \RNumPlainAddTwoRunsDone{} of the \RNumPlainAddTwoRunsTotal{} one-shard runs had finished, before anyone looked at them and before any eight-shard run.
It added eight-shard placement, live extra search, and the remaining tests, which is why one-shard runs have only planned extra search.

We call the tests in the plan registered tests; a test that one method beats another passes only if the whole range is above zero.
Under starting-center placement, we tested that refresh all gains less than \RNumPlainGateThreshold{} points on average over never refresh.
We also tested that planned extra search, live extra search, and push each beat the best refresh rule.
Under one-shard placement, we predicted that all five refresh rules match refresh all exactly.
We also tested that refresh beats extra search and push at every budget.
The probe test asked whether the probe beats each other rule by at least \RNumPlainGateThreshold{} points on average, with the whole range above zero.
It ran only if refresh all first gained at least \RNumPlainGateThreshold{} points on average, which we call the \RNumPlainGateThreshold{}-point threshold.
Under one-shard placement, the predicted tie meant that this test could not pass.
Eight-shard placement repeats these tests with live extra search.
It adds a prediction that the four rules other than the probe match refresh all at a budget of eight.

After reading all the main results, we wrote down a follow-up and stored it before running it.
In it, most-changed pays for its score bytes, so it refreshes as many summaries as the probe; we call this version paying most-changed.
Of its \RNumAddThreeTotalRuns{} simulations, \RNumAddThreeMThreeRuns{} used eight-shard placement.
Another \RNumAddThreeMOneRuns{} used starting-center placement at \RNumPlainQueriesHigh{} queries per step.
The last \RNumAddThreeRepRuns{} were reruns with the charge turned off, and they matched the stored main runs on every query.
Because we chose the follow-up after seeing the results, it is exploratory and has no pass or fail tests.
It has \RNumAddThreeTotalIntervals{} ranges, one for each of \RNumAddThreeCompRows{} comparisons on recent and on overall recall, not widened for many comparisons.
Table~\ref{tab:main} shows the eight-shard ones against the probe, push, and live extra search, on recent recall.

The main experiments have \RNumTotalMainRuns{} runs, with none missing, and every run stayed within its budget for the costs we charge.
Runs that ran out of memory were rerun from the same code.
They include some starting-center runs, two eight-shard runs, and \RNumAddThreeMThreeFailedAttempts{} of the \RNumAddThreeMThreeRuns{} eight-shard follow-up runs.
